\documentclass[12pt,onecolumn]{article}

\usepackage[utf8]{inputenc}
\usepackage[T1]{fontenc}
\usepackage{lmodern}
\usepackage[margin=1in]{geometry}
\usepackage{setspace}
\usepackage{microtype}
\usepackage{amsmath,amssymb}
\usepackage[numbers,sort&compress]{natbib}
\usepackage[colorlinks=true,linkcolor=blue,citecolor=blue,urlcolor=blue]{hyperref}
\usepackage{orcidlink}
\usepackage{fancyhdr}

\setstretch{1.10}
\setlength{\headheight}{15pt}

\pagestyle{fancy}
\fancyhf{}
\lhead{The Past Does Not Persist}
\rhead{John C. W. McKinley}
\cfoot{\thepage}

\title{
\textbf{The Past Does Not Persist,}\\
\textbf{but Its Consequences Do}\\[0.4em]
\large Dynamical State, Ongoing Change,
and the Sufficiency of the Present
}

\author{
John C. W. McKinley\,\orcidlink{0009-0005-7097-5035}\\
\small Independent Researcher
}

\date{October 2026}

\begin{document}
\maketitle

\noindent\footnotesize\textit{
Working draft published at \href{https://johncwmckinley.com}{johncwmckinley.com}.
Not deposited on Zenodo; subject to revision.
}\normalsize

\begin{abstract}

The proposition that only the present physically exists raises a foundational question: if earlier physical configurations no longer exist, how can momentum persist, motion continue, or a physical change remain underway? A present understood merely as an instantaneous arrangement of objects appears insufficient to describe a functioning physical universe.

This paper distinguishes static configuration from dynamical state. Established physical theories represent systems through more than their instantaneous positions. Momentum, angular momentum, field configurations, excitation states, and quantum information participate in determining or constraining subsequent physical behavior. These quantities belong to the physical state under consideration; they are not themselves collections of independently existing earlier configurations.

The paper further distinguishes historical records from dynamical consequences. Records preserve information concerning earlier interactions. Dynamical conditions participate in subsequent interactions. A single present physical condition may perform both functions.

Within the Timeless Light Model (TLM), these distinctions support the proposition that earlier events need not remain independently in existence for their consequences to remain physically effective. However, identifying momentum or other dynamical quantities as present-state variables does not explain their ultimate physical instantiation. Nor does it explain how one actual condition gives way to another when the governing laws are understood as timeless.

The paper explicitly identifies these questions as an unresolved state-realization problem. It distinguishes the mathematical sufficiency of a properly specified dynamical state from an ontological explanation of why that state is actual and how lawful succession occurs.

The result is a clarification rather than a completed physical mechanism. The past need not persist as an independently existing physical reality for its consequences to remain effective. But the dynamical content and lawful continuation of the present remain subjects requiring further foundational explanation.

\end{abstract}


\section{Introduction: The Baseball in Flight}

Consider a baseball traveling through the air.

At the physical condition under consideration, the ball occupies a particular position relative to the Earth, possesses momentum, and participates in gravitational and electromagnetic interactions with its surroundings.

The ordinary description is that the baseball has traveled along a trajectory and will continue moving until its motion is altered.

That description is useful and predictive. It permits calculations of the ball's future position, the likely point of impact, and the physical consequences of its eventual collision.

But a different question arises when considering the ontological status of the past.

Must the earlier positions of the baseball continue physically existing in order for the ball to possess its present momentum?

If the earlier configurations no longer exist, what distinguishes a baseball moving upward from one moving downward when both occupy the same position?

What, precisely, makes the ball's motion ongoing?

These questions reveal an important distinction between a physical configuration and a complete dynamical state.

A description consisting only of the ball's position, mass, and visible arrangement is incomplete. It does not specify the information needed to describe the ball's subsequent behavior.

Classical mechanics includes momentum among the variables specifying the physical state.

The earlier trajectory need not be retained as a collection of physically existing positions in order for the present state to include momentum.

This observation appears to resolve an immediate difficulty for a present-only ontology.

Yet a deeper question remains.

If momentum is a present physical property, what accounts for its physical reality? What makes that property part of the actual state rather than merely a number assigned in a mathematical description?

And what connects the actual state under consideration to the next actual physical condition?

The present paper separates these questions.

It argues that the continued physical existence of the past is unnecessary for the standard dynamical description of a present state.

It does not claim that naming momentum as a present property explains the ultimate physical mechanism of dynamical succession.

This distinction is essential to the Timeless Light Model (TLM), which proposes singular physical actuality governed by timeless law.


\section{The Past as History and the Past as Existence}

The statement that the past does not persist requires careful definition.

It does not mean that earlier events never occurred.

A collision may have occurred. Momentum may have been exchanged. Energy may have been transferred. Physical systems may have been deformed, heated, accelerated, or otherwise altered.

The resulting conditions may remain physically present after the original interaction has ended.

What the proposition denies is the independent continued existence of the earlier configuration as a separately actual physical condition.

An event having occurred and an event continuing to exist are different claims.

Consider a hammer striking a metal plate.

The impact occurs and the plate is deformed.

The deformation persists as a physical condition after the impact has ended.

The plate does not need an independently existing copy of its earlier undeformed configuration for the deformation to remain real.

The present plate possesses its resulting physical organization.

The same distinction applies when the consequence is dynamical rather than a permanent deformation.

A collision may impart momentum to a body.

Afterward, the body possesses a momentum state associated with its subsequent physical behavior.

The earlier collision is not a force continuously acting upon the body from an independently existing past.

The momentum is part of the resulting physical state.

Within TLM, the proposed ontology is that physical actuality is singular. Earlier configurations do not remain separately actual merely because they can be remembered, represented mathematically, or reconstructed from physical evidence.

The present argument asks whether the resulting physical actuality includes the information required for continuing physical behavior.

The answer supplied by established dynamical descriptions is that configuration alone is insufficient.

The state must include additional dynamical information.

Whether that mathematical distinction also supplies a complete ontology of physical succession is a separate question.


\section{Historical Records and Dynamical Consequences}

The consequences of earlier interactions have at least two distinguishable functions.

One is historical.

The other is dynamical.

A historical record contains information from which an earlier physical situation may be reconstructed.

A photograph, electronic memory, geological formation, deformed material, or accumulated clock reading provides a familiar example.

Each is a presently existing condition.

Its informational significance concerns earlier interactions.

A dynamical consequence, by contrast, is a presently relevant condition that participates in determining or constraining subsequent physical behavior.

Momentum, angular momentum, field configurations, stored energy, and excitation states provide examples.

Consider a baseball photographed in flight.

The photograph preserves information about the ball's appearance and position during the exposure.

The ball's momentum, however, is part of its dynamical description.

The photograph is a record of an interaction.

The momentum participates in the description of subsequent interactions.

Both belong to present physical reality, but they perform different explanatory functions.

The distinction is functional, not a division into mutually exclusive physical substances.

A compressed spring may preserve evidence of an earlier applied force while also possessing potential energy relevant to its subsequent behavior.

A rotating flywheel may preserve a measurable consequence of an earlier torque while its current angular momentum governs its response to further torques.

The same physical condition therefore may serve both historical and dynamical purposes.

The distinction concerns the question being asked.

Historical reconstruction asks what earlier physical circumstances produced the evidence.

Dynamical analysis asks what behavior is admissible from the specified physical state.

The past does not need to remain physically in existence for either function to be meaningful.

However, the second function introduces a further ontological question.

If a dynamical property belongs to the present, what is its actual physical status?

The question cannot be answered simply by relabeling the property as a record.

Momentum is not merely a memory of motion.

Its role in present physical dynamics requires a more precise account.


\section{Configuration Is Not Dynamical State}

Suppose two otherwise comparable hypothetical situations place a baseball at the same position above the ground.

In one situation, the ball possesses upward momentum.

In the other, it possesses downward momentum.

The positions are identical.

The dynamical states are different.

A description consisting only of the present position cannot distinguish the two situations.

Classical Hamiltonian mechanics provides a clear mathematical illustration \cite{Goldstein2002}.

For a system described by generalized coordinates $q_i$ and conjugate momenta $p_i$, the state may be represented as

\begin{equation}
S=(q_1,\ldots,q_n,p_1,\ldots,p_n).
\end{equation}

The coordinates represent configuration variables.

The conjugate momenta supply additional dynamical information.

For a Hamiltonian $H(q,p,t)$, the equations of motion are

\begin{equation}
\dot q_i=\frac{\partial H}{\partial p_i},
\qquad
\dot p_i=-\frac{\partial H}{\partial q_i}.
\end{equation}

For a simple nonrelativistic particle in a potential $V(q)$,

\begin{equation}
H(q,p)=\frac{p^2}{2m}+V(q),
\end{equation}

giving

\begin{equation}
\dot q=\frac{p}{m},
\qquad
\dot p=-\frac{\partial V}{\partial q}.
\end{equation}

These equations distinguish a specification of position from a specification of dynamical state.

The difference is important.

A position describes where a body is within the chosen coordinate representation.

Momentum contributes information relevant to its lawful motion and interactions.

For an adequately specified classical system, the state and governing dynamics supply the mathematical basis for predicting subsequent behavior.

The equations do not require a separate physical inventory of every earlier state.

They also do not prove that earlier states have ceased to exist.

The mathematical formalism is compatible with different ontological interpretations.

Its contribution to the present argument is the distinction between the information contained in a state and the historical sequence from which that state may be inferred or reconstructed.

The present state need not contain a complete historical record in order to provide the information required by its dynamical equations.

This observation helps clarify the proposed sufficiency of the present.

It does not yet explain why the actual universe possesses one particular dynamical state rather than another.


\section{Must the Present Remember Where It Was?}

The baseball example permits the central question to be stated more directly.

Suppose a physical system must possess some present information indicating its direction and speed.

Does the system need a record of where it was previously?

In classical Hamiltonian mechanics, the answer is no.

A properly specified state includes the momentum variables required by the equations of motion.

The equations do not demand that the system first reconstruct an earlier sequence of positions before assigning its present velocity.

The state variables already contain the necessary dynamical information.

That does not mean that all physical models are memoryless in their usual mathematical formulation.

Some descriptions of selected subsystems include memory terms or explicit dependence upon earlier conditions.

Such descriptions arise, for example, when relevant environmental degrees of freedom have been excluded from the explicit state representation.

A broader description may include additional environmental variables and correlations.

The important distinction is between a physically complete description and an intentionally reduced one.

A model may require historical data because its selected variables do not fully describe the relevant physical situation.

That mathematical dependence does not establish that earlier physical configurations continue to exist independently.

The proposed TLM distinction is therefore not that every practical equation is memoryless.

It is that a physical history represented in an equation is not automatically an independently existing physical object.

For a state-based dynamical theory, the current state supplies the information specified by the theory.

A system does not need to retain an indefinitely extending inventory of earlier locations merely to possess a present momentum.

The unresolved question is not whether momentum appears in the mathematical state.

It does.

The unresolved question is what physically constitutes that dynamical information within the proposed ontology.


\section{Momentum Is Not a Historical Record}

In elementary nonrelativistic mechanics,

\begin{equation}
p=mv,
\end{equation}

where velocity is commonly expressed as

\begin{equation}
v=\frac{dx}{dt}.
\end{equation}

The derivative describes a local mathematical relationship within a represented trajectory.

It does not require every point of the trajectory to exist simultaneously.

A trajectory is a mathematical description of lawful relationships among physical states.

Its use does not settle the ontological status of the past.

Momentum may be measured or inferred through observations made over an interval.

That measurement procedure does not establish that all the earlier configurations involved in the procedure remain physically actual.

The distinction is between the method by which a physical quantity is determined and the role the quantity plays in the dynamical description.

A present property is not identical to the historical procedure used to measure or infer it.

In an appropriate classical state description, momentum is a dynamical variable.

Its mathematical role does not require an independently retained past trajectory.

Nevertheless, one must avoid turning the mathematical variable into a complete physical explanation.

Writing the symbol $p$ supplies a representation of momentum.

It does not, by itself, explain why a particular physical system possesses that momentum.

The present paper therefore makes a narrower claim than a complete physical account of motion.

It claims that a presently specified momentum does not logically require the continued existence of every earlier position associated with the body's history.

It does not claim that the ontological nature of momentum has thereby been fully explained.


\section{Ongoing Change as a Present Condition}

The expression \emph{in progress} ordinarily refers to a process that began earlier and has not yet reached its completion.

A chemical reaction is in progress.

A flywheel is rotating.

A projectile is moving toward a target.

An excited atom has not yet undergone a particular transition.

Within a present-only ontology, these descriptions require careful interpretation.

An ongoing change is not a physical fragment of the past that remains active behind the present.

It is represented by presently relevant dynamical conditions and their lawful relationships.

A chemical reaction involves actual molecular species, energies, concentrations, fields, and other conditions.

These participate in governing the admissible subsequent reaction pathways.

A rotating flywheel possesses angular momentum and physical relationships that govern its response to applied torques.

An excited atom possesses a quantum state relevant to its admissible transitions and interactions.

These conditions are not independently existing earlier events.

They belong to the physical state being described.

The expression \emph{in progress} therefore need not imply that a portion of the past remains physically in existence.

It identifies a present physical situation associated with further lawful behavior.

There is also no need to postulate a separate metaphysical label reading \emph{in progress}.

The phrase summarizes the dynamical relationships already represented by the relevant physical state.

That distinction matters because adding an unexplained status label would merely rename the original problem.

The physical description must identify the relevant state variables and laws.

The ontology must then address what makes those variables and relationships physically actual.

The first task belongs to dynamical modeling.

The second is the foundational question under examination.


\section{The Unresolved State-Realization Problem}

The preceding analysis clarifies what an adequate dynamical description requires.

It does not yet explain the ultimate physical realization of the state.

This is a genuine unresolved issue within the proposed TLM ontology.

Consider again two hypothetical physical situations with the same spatial configuration but different momenta.

The mathematical description distinguishes them by assigning different state variables.

But why is one momentum physically actual rather than another?

What, in the ontology, constitutes the physical difference between them?

It is insufficient to answer that momentum is present because the present contains momentum.

That statement identifies the required property without explaining its physical realization.

Nor does appealing to a previous trajectory resolve the foundational question if the theory has already rejected the continued independent existence of that trajectory.

The earlier trajectory may explain how the present condition arose.

It does not establish what presently makes the dynamical condition physically actual.

At least three questions must therefore be distinguished.

First, there is the \emph{state specification} question.

What information must an adequate physical description contain?

Established dynamical theories provide answers through their state variables, constraints, and equations.

Second, there is the \emph{state realization} question.

What makes a particular specified state the actual physical condition rather than merely one mathematically admissible possibility?

The identification of state variables does not by itself answer this.

Third, there is the \emph{succession} question.

What accounts for one actual physical condition giving way to another when the underlying governing laws are regarded as timeless?

The mathematical equations describe relationships among states.

They do not, merely by being written, establish a unique ontology of physical actualization.

These questions define the remaining black box.

The present proposal does not conceal it.

Instead, it isolates the unresolved issue from a different and already clarified one.

An existing past is not needed as an additional collection of physical configurations for a state-based mathematical description.

But a complete present-only ontology must still explain what physical actuality consists of and how lawful dynamical succession is realized.

The distinction is essential to evaluating the completeness of TLM.


\section{Timeless Laws and Physical Succession}

TLM proposes that the fundamental architecture of reality is the totality of timeless governing laws.

These laws are not themselves objects situated in space or processes accumulating elapsed time.

They govern the admissible relationships, interactions, and outcomes of physical actuality.

A physical state, by contrast, specifies the actual conditions relevant to the interaction under consideration.

The distinction may be represented schematically as

\begin{equation}
\boxed{
\left(\mathcal{L},S_{\mathrm{actual}}\right)
\longrightarrow \mathcal{O}_{\mathrm{admissible}}
}
\end{equation}

where $\mathcal{L}$ denotes the timeless governing structure, $S_{\mathrm{actual}}$ denotes the actual physical state, and $\mathcal{O}_{\mathrm{admissible}}$ denotes the admissible outcomes or interactions.

This expression is a conceptual schema.

It is not a new equation of motion.

It also does not identify the mechanism by which an admissible outcome becomes actual.

A deterministic dynamical theory may assign a unique subsequent state to sufficiently specified conditions.

A probabilistic theory may assign an outcome distribution.

Neither mathematical description alone determines the ontological meaning of actualization.

A schematic succession relation might be written as

\begin{equation}
S_{n+1}=F_{\mathcal{L}}(S_n),
\end{equation}

for a deterministic discrete description.

But the notation must not be mistaken for a derivation of physical succession.

The index $n$ orders mathematical states.

It does not establish that the universe possesses a universal sequence of discrete ticks.

Nor does the function $F_{\mathcal{L}}$ explain why the represented next state becomes actual.

In standard dynamical theories, the equations express lawful relationships involving time or other evolution parameters.

TLM's further ontological question is how those relationships are to be understood when the governing lawful structure is timeless.

The existence of a rule does not by itself supply a separate account of its physical instantiation.

Likewise, the existence of a mathematical state description does not automatically explain why that state is actual.

A complete theory must distinguish these categories without introducing an unexplained temporal mechanism under a different name.

The present paper does not claim to have completed that derivation.

Its contribution is to specify more clearly the explanatory work remaining.


\section{Possibility Does Not Imply Multiple Actual Universes}

The incompleteness of a static configuration suggests another important question.

If the same position is compatible with many possible momentum states, does every possible state require its own physically existing universe?

The mathematical answer does not require that conclusion.

Consider the same position $q$ combined with different possible momenta:

\begin{equation}
S_1=(q,p_1),\qquad
S_2=(q,p_2),\qquad
S_3=(q,p_3).
\end{equation}

The state space may admit a continuous range of possibilities.

Those possibilities are mathematically distinguishable.

Their admissibility does not establish that all are simultaneously actual.

The distinction between possible and actual states is therefore fundamental.

Within TLM, the timeless governing structure constrains admissible physical relationships and outcomes.

The physical actuality under consideration is not identified with the set of every mathematically admissible state.

The existence of many possible momentum values does not, by itself, require many actual universes.

This observation does not resolve the quantum measurement problem or exclude every interpretation involving multiple branches or worlds.

It simply distinguishes mathematical possibility from ontological multiplicity.

The outstanding TLM question remains:

What establishes the particular actual physical condition within the admissible structure?

The mere number of possibilities does not answer it.


\section{Quantum States and Dynamical Information}

The classical variables of position and momentum do not exhaust the dynamical information used in quantum physics.

Quantum mechanics requires an appropriate quantum state description \cite{Sakurai2017}.

For a closed system, Schr\"odinger evolution is expressed as

\begin{equation}
i\hbar\frac{\partial}{\partial t}|\psi(t)\rangle
=\hat H|\psi(t)\rangle.
\end{equation}

For a density operator $\rho$, unitary evolution is represented by

\begin{equation}
i\hbar\frac{d\rho}{dt}=[\hat H,\rho].
\end{equation}

These equations provide the quantum dynamical relationships appropriate to the stated conditions.

They do not generally assign simultaneous definite classical positions and momenta to a quantum system.

The present argument does not require them to do so.

Its central point is that whatever state information the applicable physical theory requires must not be confused with a collection of independently existing earlier configurations.

For a projective measurement described by projectors $P_i$, the Born rule gives

\begin{equation}
p_i=\operatorname{Tr}(\rho P_i).
\end{equation}

The state and measurement description specify the probabilities of the permitted outcomes.

The physical registration of a particular outcome remains a further question.

In TLM terminology, the distinction is between lawful eligibility and actual physical registration.

That distinction is compatible with the proposed separation between timeless law and physical actuality.

But it also illustrates the state-realization problem.

The mathematical description of admissible outcomes does not, by itself, explain why one particular result is physically registered.

Likewise, describing a quantum state does not automatically settle the ultimate ontology of that state.

The present paper therefore does not claim that the dynamical completeness of a quantum state solves the measurement problem.

Its narrower claim concerns the distinction between present state information and independently existing physical history.

That distinction remains meaningful without imposing a classical trajectory ontology on quantum systems.


\section{Relativity and the Meaning of the Present}

The expression \emph{universal present} requires a distinction between physical actuality and coordinate simultaneity.

Special relativity does not provide a preferred universal assignment of simultaneous distant events \cite{Einstein1905,TaylorWheeler1992}.

For a Lorentz transformation along one spatial direction,

\begin{equation}
t'=\gamma\left(t-\frac{vx}{c^2}\right),
\end{equation}

where

\begin{equation}
\gamma=\frac{1}{\sqrt{1-v^2/c^2}}.
\end{equation}

Events assigned the same coordinate time in one inertial frame need not be assigned the same coordinate time in another.

This is a structural feature of relativistic spacetime.

It cannot be removed merely by compensating for signal-travel delay.

TLM's proposition of singular physical actuality is therefore not identified here with a preferred global simultaneity coordinate.

The present paper does not introduce a universal master clock or a physically preferred foliation of spacetime.

In relativity, dynamical conditions must be expressed through appropriate local and covariant quantities.

For a massive particle,

\begin{equation}
p^\mu=mu^\mu,
\end{equation}

with

\begin{equation}
u^\mu=\frac{dx^\mu}{d\tau},
\end{equation}

where $\tau$ is proper time.

For field and gravitational systems, the relevant description includes the appropriate field variables, stress-energy, geometry, and dynamical relationships \cite{Wald1984}.

Suitable relativistic field theories may admit initial-value formulations using data specified on spacelike hypersurfaces.

That mathematical procedure does not establish a uniquely privileged universal physical present.

The present paper's claim is therefore not that every distant physical condition has been assigned to a new absolute time slice.

Its claim is that physical actuality, however ultimately represented, must not be reduced to static position information.

The ontological interpretation must preserve the objective dynamical relationships and causal constraints of relativity.

In particular, the registration of information remains subject to relativistic causal structure.

Delayed registration is not equivalent to delayed existence.

Local proper time is not a universal clock.

And the mathematical representation of a dynamical state does not settle the deeper problem of how singular physical actuality is to be formulated relativistically.


\section{Connection to Space Is Not a Box}

The related paper \emph{Space Is Not a Box} developed a relational interpretation of spatial experience, motion, and physical registration within TLM \cite{McKinleySpaceBox2026}.

Its central argument was that the success of spatial descriptions does not require an independently existing spatial container.

The paper also distinguished physical actuality from the delayed availability of physical information and the present records of earlier events.

The present work extends that analysis by addressing the dynamical content of actuality.

The earlier argument explained how lawful relationships and registered consequences support the reconstruction of motion and physical history.

The present paper emphasizes that physical motion is not exhausted by its historical reconstruction.

The meteor crossing the night sky provides an example.

Observers register changes in apparent position, brightness, and direction.

Their records support the reconstruction of a trajectory.

That reconstruction has objective physical content.

But the meteor's physical condition is not exhausted by the reconstructed history.

Its momentum, thermal state, interaction with the atmosphere, and surrounding field conditions participate in governing its subsequent behavior.

Those dynamical conditions belong to the physical description.

They do not require a surviving physical trail of earlier meteor configurations.

Thus the two papers address complementary questions.

\emph{Space Is Not a Box} asks whether physical relationships require an independent spatial container.

The present paper asks whether continuing physical behavior requires an independently existing past.

The proposed answer is that the present dynamical state and governing laws provide the relevant physical description without adding such a past.

The new question identified here concerns the physical reality of that state.

The elimination of an independently existing past does not eliminate the need to explain what makes present dynamical conditions actual.

This is a substantive foundational requirement rather than a matter of terminology.


\section{What the Model Establishes and What Remains Unexplained}

The logical status of the argument must be stated explicitly.

Established physical theories distinguish configuration from the further information required for dynamical description.

Classical Hamiltonian mechanics provides one example.

Relativistic field theories and quantum mechanics provide different and more general state structures.

These theories do not require earlier configurations to remain physically in existence as additional inputs to every state-based dynamical calculation.

This mathematical observation supports the coherence of distinguishing a present state from its historical reconstruction.

It does not prove that the past has no independent ontological existence.

That proposition belongs to TLM's foundational interpretation.

The present paper also does not derive a new physical law, introduce a new force, or alter established quantitative predictions.

Most importantly, it does not claim to have explained the physical mechanism by which a present state possesses its dynamical properties.

Writing

\begin{equation}
S=(q,p)
\end{equation}

identifies variables in a classical state description.

It does not, by itself, explain how momentum is physically instantiated.

Writing an equation of motion identifies a lawful relationship.

It does not, by itself, supply an ontology of actualization.

And asserting that only one physical actuality exists does not, by itself, explain how its dynamical conditions are realized or how succession is obtained.

These distinctions mark the present boundary of the argument.

The remaining theoretical task is to connect timeless governing law, actual dynamical state, and lawful physical succession without merely reintroducing an unexplained temporal container or an existing past under another name.

That task is not resolved here.

Recognizing it prevents a description of dynamical state from being mistaken for a completed explanation of dynamical reality.


\section{Discussion: What the Present Actually Carries}

The ordinary description of the baseball begins with a history.

The ball was struck, acquired momentum, and traveled through a sequence of positions.

That history may be correct.

But the ball's present physical condition is not identical to a surviving collection of those earlier positions.

Its current dynamical state contains information relevant to its lawful interactions.

The distinction between historical explanation and dynamical description is central.

The question of how the ball acquired momentum concerns its physical history.

The question of how the ball behaves from its specified current state concerns its dynamics.

These are related but different questions.

The first does not require the historical event to continue existing.

The second does not require a complete historical record as an additional physical ingredient when the relevant dynamical state has been adequately specified.

That is the conceptual resolution advanced here.

Yet the state itself remains physically significant.

If two otherwise identical present configurations possess different momentum states, something about their physical conditions differs.

A complete ontology must account for that difference.

Calling it momentum supplies a precise physical quantity.

It does not finish the ontological explanation.

There is therefore an important distinction between two meanings of sufficiency.

The present may be sufficient for a physical theory's dynamical description when the required state variables, constraints, and governing laws are supplied.

But the mere specification of those variables is not necessarily a sufficient explanation of why that particular state is actual.

The first is a mathematical claim about state descriptions.

The second is a foundational claim about physical reality.

The present paper establishes their separation.

Within TLM, the governing laws are timeless.

Physical actuality is singular.

Its actual conditions include the consequences of earlier interactions.

Records permit historical reconstruction.

Dynamical conditions participate in ongoing physical behavior.

The independent continued existence of earlier configurations is not required by those descriptions.

But the realization of the dynamical state and the actualization of subsequent conditions remain theoretical questions.

The model must ultimately account for both.

The present is not a frozen snapshot.

Neither is the word \emph{present} an explanation of physical dynamics.

The content of physical actuality must be specified, and its lawful succession must be understood.


\section{Conclusion}

The proposition that only the present physically exists does not imply that reality is static.

Physical state descriptions contain more than positions and arrangements.

Momentum, angular momentum, field configurations, excitation states, and quantum information participate in the lawful description of physical behavior.

These dynamical conditions are not identical to independently existing historical trajectories.

Earlier events may have occurred and may have produced lasting physical consequences.

The earlier configurations do not need to remain independently in existence for those consequences to be physically relevant.

The present therefore contains both historical records and dynamical conditions, whose explanatory functions are distinguishable.

Within TLM, this observation supports an ontology in which physical actuality is singular and the governing lawful structure is timeless.

It also identifies an unresolved foundational issue.

A mathematical state variable is not, by itself, an explanation of its physical realization.

A mathematical law of evolution is not, by itself, an explanation of why one actual condition gives way to another.

A completed TLM account must ultimately address these questions without requiring an independently existing past as an additional physical substrate.

The contribution of the present paper is therefore twofold.

It clarifies why the continued existence of earlier physical configurations is unnecessary for the standard dynamical description of a properly specified present state.

And it identifies the remaining ontological problem of state realization and lawful succession.

The past need not persist for its consequences to remain effective.

But the dynamical reality of the present still requires explanation.

The past does not persist.

Its consequences do.


\begin{thebibliography}{99}

\bibitem{Goldstein2002}
H.~Goldstein, C.~Poole, and J.~Safko.
\newblock \emph{Classical Mechanics}.
\newblock 3rd edition,
Addison-Wesley (2002).

\bibitem{Einstein1905}
A.~Einstein.
\newblock Zur Elektrodynamik bewegter K{\"o}rper.
\newblock \emph{Annalen der Physik},
\textbf{17}, 891--921 (1905).

\bibitem{TaylorWheeler1992}
E.~F.~Taylor and J.~A.~Wheeler.
\newblock \emph{Spacetime Physics}.
\newblock 2nd edition,
W.~H.~Freeman (1992).

\bibitem{Wald1984}
R.~M.~Wald.
\newblock \emph{General Relativity}.
\newblock University of Chicago Press (1984).

\bibitem{Sakurai2017}
J.~J.~Sakurai and J.~Napolitano.
\newblock \emph{Modern Quantum Mechanics}.
\newblock 2nd edition,
Cambridge University Press (2017).

\bibitem{McKinleySpaceBox2026}
J.~C.~W.~McKinley.
\newblock \emph{Space Is Not a Box:
Timeless Law, Universal Actuality,
and the Registration of Space and Motion}.
\newblock Working draft, October 2026.
\newblock \url{https://johncwmckinley.com}.

\end{thebibliography}

\end{document}
